\section{A uniform presentation for a tree}
Let $T$ be a finite tree on $n$ vertices. Write $s_e$ for its edge
transpositions. The following is the presentation in Fernandes--Leemans,
Proposition 3.1, without its unnecessary restriction $n\geq7$.
\begin{theorem}\label{thm:tree}
The symmetric group on the vertices of $T$ has generators $s_e$ and
relations
\[
 s_e^2=1,\qquad (s_es_f)^2=1\quad(e,f\text{ disjoint}),\qquad
 (s_es_f)^3=1\quad(e,f\text{ meet}),
\]
together with $(s_es_fs_es_g)^2=1$ for every three distinct edges
$e,f,g$ incident with one vertex. For each unordered triple it is enough to impose the relation in one
ordering; all the other orderings then follow.
\end{theorem}
\begin{proof}
First observe that the choice of ordering is immaterial. For three
incident involutions $a,f,g$, the pairwise braid relations give
\[
 [afa,g]=1\ \Longleftrightarrow\ [a,fgf]=1
 \ \Longleftrightarrow\ [a,gfg]=1
 \ \Longleftrightarrow\ [aga,f]=1.
\]
The first and last equivalences follow by conjugation and a braid
identity. Also $afa=faf$ interchanges the first two letters.
These interchanges generate all orderings of the triple.

Let $P(T)$ be the presented group. The indicated transpositions satisfy
the relations and generate $\Sym(V(T))$, so there is a surjection
$P(T)\to\Sym(V(T))$. We prove $|P(T)|\leq n!$ by induction on $n$.
The cases $n\leq2$ are immediate. Remove a leaf $p$, with neighbor $q$,
and write $a=s_{pq}$. By induction the subgroup $H$ generated by the
remaining edge generators is a quotient of $\Sym(V(T)\setminus\{p\})$.
Its displayed permutation image is that entire symmetric group, so the
quotient is faithful and we may identify $H$ with $S_{n-1}$.

We first prove that $a$ centralizes $H_q$, the point stabilizer of $q$
in $H$. It commutes with every edge not incident with $q$. If
$f=s_{qu}$ and $g=s_{qv}$ are two distinct other edges at $q$, the
triangle relation says that $afa$ commutes with $g$. The braid
identity $afa=faf$ turns this into the assertion that $a$ commutes
with $fgf=(uv)$. The edge transpositions within the components of
$T-\{p,q\}$, together with these transpositions between their neighbors
at $q$, form a connected transposition graph on all vertices other than
$p,q$. They therefore generate $H_q$. This proves the assertion,
including the case where there is at most one such component.

Choose a neighbor $u$ of $q$ in $T-p$. The relation for $a$ and
$s_{qu}$ is the braid relation. Conjugating it by $H_q$, which acts
transitively on the remaining $n-2$ vertices and centralizes $a$, gives
\[
 af_va=f_vaf_v\qquad\text{for }f_v=(qv),\quad v\ne p,q.
\]
Set $f_q=1$. Every element of $H$ has a unique left coset of $H_q$,
represented by some $f_v$. Consequently the union of the following
at most $n$ left cosets of $H$ is stable under right multiplication by
$H$:
\[
 H,\qquad Ha f_v\quad(v\ne p).
\]
Explicitly, $f_vh=kf_w$ for some $k\in H_q$, so
$Ha f_vh=Ha kf_w=Ha f_w$. It is stable under multiplication by $a$
as well: $Ha\cdot a=H$, and for $v\ne q$ the braid identity gives
$Ha f_va=Hf_va f_v=Ha f_v$.
The union contains the identity, and $P(T)$ is generated by $H$ and
$a$. It is therefore all of $P(T)$, whence $[P(T):H]\leq n$ and
$|P(T)|\leq n!$. The surjection to $S_n$ is an isomorphism.
\end{proof}

\section{Transporting an exact presentation}
\begin{lemma}\label{lem:transport}
Suppose $K=\gen{s_e\mid R(s)}$. Let $r_i=v_i(s)$ generate $K$ and
choose words $w_e(r)$ with $w_e(v(s))=s_e$ in $K$. Then an exact
presentation on the marking $r_i$ is
\begin{equation}\label{eq:transport}
 \gen{r_i\mid R(w(r)),\quad r_i=v_i(w(r))\text{ for every }i}.
                                                               \end{equation}
\end{lemma}
\begin{proof}
The prescribed elements of $K$ satisfy all displayed relations, giving
a surjective map from \eqref{eq:transport} to $K$. Conversely $s_e\mapsto w_e(r)$
defines a homomorphism from $K$ to \eqref{eq:transport}, because every relator $R(w)$
is imposed. One composite fixes each $s_e$ by the word identity in
$K$; the other fixes each $r_i$ by the second group of relations.
Thus these homomorphisms are inverse.
\end{proof}

